\SourceLeaf{357}
\par\bigskip\Needspace{7\baselineskip}
\pdfbookmark[1]{On the Theory of Biquadratic Residues. V.}{nachlass-biquadratic-V}
\centerline{[V.]}\medskip
\renewcommand{\Article}[1]{\par\medskip\Needspace{4\baselineskip}\pdfbookmark[2]{Art. #1}{nachlass-biquadratic-V-art-#1}\centerline{\textbf{[#1.]}}\nobreak\smallskip}
\Article{1}
\[
\begin{array}{lll}
\text{Modulus}&M=A+Bi&AA+BB=D\\
\text{Residue}&m=a+bi,&aa+bb=d
\end{array}
\]
\[
\frac{mD}{M}=\mu=\alpha+\beta i=aA+bB+(Ab-Ba)i
\]
\[
\xi+\eta i=\pi,\quad\pi m=p=x+yi;\quad\pi M=P=X+Yi
\]
$\omega$ is an undetermined infinitely small real positive quantity.

\Article{2}
\begin{center}\emph{Preparation.}\end{center}
I. Collect all $\pi$ for which
\[
\begin{array}{l}
\xi\text{ nonnegative and no greater than }\frac12\\
\eta\text{ positive and less than }\frac12\\
\text{Either }x\text{ or }y\text{ is an integer}\\
\text{Either }X\text{ or }Y\text{ is an integer}
\end{array}
\]
and, for each $\pi$, determine the quantity $\varepsilon$ by the following rule:
\[
\begin{array}{ll}
\text{Let }p^0&\text{be the nearest integer divisible by }1+i\text{ to }p\\
P^0&\text{be the nearest integer divisible by }1+i\text{ to }P
\end{array}
\]
and put $\varepsilon=\pm1$, where the sign is always the same as the sign of the imaginary part of the quantity
\[
\frac{p-p^0}{P-P^0}(\alpha-\beta i)
\]

The following are the special rules. First class, $x$ and $X$ integers
\[
\begin{aligned}
\beta\xi&=-Bx+bX\\
\beta\eta&=-Ax+aX\\
\beta y&=-\alpha x+dX\\
\beta Y&=-Dx+\alpha X
\end{aligned}
\]

\SourceLeaf{358}
$\varepsilon=-1$ when $\beta$ is positive, $x$ even, $[y]$ even, $X$ even, $[Y]$ even, or when only an odd number of these conditions holds.

$\varepsilon=+1$ when none or an even number of these conditions holds.

Second class, $y$ and $Y$ integers
\[
\begin{aligned}
\beta\xi&=+Ay-aY\\
\beta\eta&=-By+bY\\
\beta x&=+\alpha y-dY\\
\beta X&=+Dy-\alpha Y
\end{aligned}
\]
$\varepsilon=-1$ when $\beta$ is positive, $[x]$ even, $y$ even, $[X]$ even, $Y$ even, or when an odd number of these conditions holds.

$\varepsilon=+1$ when none or an even number holds.

Third class, $y$ and $X$ integers
\[
\begin{aligned}
\alpha\xi&=+By+aX\\
\alpha\eta&=+Ay-bX\\
\alpha x&=-\beta y+dX\\
\alpha Y&=+Dy-\beta X
\end{aligned}
\]
$\varepsilon=-1$ when $\alpha$ is positive, $[x],y,X,[Y]$ are all even, or when an odd number of these conditions holds.

$\varepsilon=+1$ when none or an even number holds.

Fourth class, $x$ and $Y$ integers
\[
\begin{aligned}
\alpha\xi&=+Ax+bY\\
\alpha\eta&=-Bx+aY\\
\alpha y&=+\beta x+dY\\
\alpha X&=+Dx+\beta Y
\end{aligned}
\]
$\varepsilon=+1$ when $\alpha$ is positive, $x,[y],[X],Y$ are all even, or when an odd number of these conditions holds.

$\varepsilon=-1$ when none or an even number of them holds.

\SourceLeaf{359}
\Needspace{22\baselineskip}
\Article{3}
II. Collect all $\pi$ for which
\[
\begin{array}{l}
\xi\text{ positive and less than }\frac12\\
\eta=\omega\\
\text{and either }X\text{ or }Y\text{ is an integer,}
\end{array}
\]
and put $\varepsilon=\pm1$, taking the sign of the imaginary part of
\[
\frac{M}{P-P^0}
\]

Special rule: first class, $X$ an integer
\[
\begin{aligned}
A\xi&=+X+B\omega\\
Ax&=+aX-\beta\omega\\
Ay&=+bX-\alpha\omega\\
AY&=+BX+D\omega
\end{aligned}
\]
$\varepsilon=-1$ when $A$ is positive, $X$ and $[Y]$ are even, or when only one condition holds.

$\varepsilon=+1$ when none or two hold.

Second class, $Y$ an integer
\[
\begin{aligned}
B\xi&=+Y-A\omega\\
Bx&=+aY-\alpha\omega\\
By&=+bY-\beta\omega\\
BX&=+AY-D\omega
\end{aligned}
\]
$\varepsilon=+1$ when $B$ is positive, $[X]$ and $Y$ are even, or when only one condition holds.

$\varepsilon=-1$ when none or two hold.

\clearpage
\Article{4}
III. Collect all $\pi$ for which
\begin{quote}
$\xi$ and $\eta$ satisfy the same conditions as in II;\par
either $x$ or $y$ is an integer,
\end{quote}
and put $\varepsilon=\pm1$ with the sign of the imaginary part of
\[
\frac{m}{p-p^0}
\]

\SourceLeaf{360}
Special rules: first class, $x$ an integer
\[
\begin{aligned}
a\xi&=+x+b\omega\\
ay&=+bx+d\omega\\
aX&=+Ax+\beta\omega\\
aY&=+Bx+\alpha\omega
\end{aligned}
\]
$\varepsilon=-1$ when $a$ is positive, $x,[y]$ are both even, or when only one condition holds.

$\varepsilon=+1$ when none or two hold.

Second class, $y$ an integer
\[
\begin{aligned}
b\xi&=+y-a\omega\\
bx&=+ay-d\omega\\
bX&=+Ay-\alpha\omega\\
bY&=+By+\beta\omega
\end{aligned}
\]
$\varepsilon=+1$ when $b$ is positive, $[x]$ and $y$ are even, or when only one condition holds.

$\varepsilon=-1$ when none holds.

\clearpage
\Article{5}
IV. Collect all $\pi$ for which
\[
\begin{array}{l}
\xi=\frac12+\frac12\omega\\
\eta\text{ positive and less than }\frac12\\
X\text{ or }Y\text{ is an integer,}
\end{array}
\]
and put $\varepsilon=\pm1$ with the sign of the imaginary part of
\[
\frac{iM}{P-P^0}
\]
Special rules: first class, $X$ an integer,
\[
\begin{aligned}
2B\eta&=+A-2X+A\omega\\
2Bx&=-\beta+2bX-\beta\omega\\
2By&=+\alpha-2aX+\alpha\omega\\
2BY&=+D-2AX+D\omega
\end{aligned}
\]
$\varepsilon=+1$ when $B$ is positive, $X,[Y]$ are even, or when one condition holds,

$\varepsilon=-1$ when none or two conditions hold.

\SourceLeaf{361}
Second class, $Y$ an integer
\[
\begin{aligned}
2A\eta&=-B+2Y-B\omega\\
2Ax&=+\alpha-2bY+\alpha\omega\\
2Ay&=+\beta+2aY+\beta\omega\\
2AX&=+D-2BY+D\omega
\end{aligned}
\]
\noindent $\varepsilon=+1$ when $A$ is positive, $[X],Y$ are even, or when one holds;\par
\noindent $\varepsilon=-1$ when none or two conditions hold.

\clearpage
\Article{6}
V. Collect all $\pi$ for which
\[
\begin{gathered}
\xi,\eta\text{ satisfy the same conditions as in IV,}\\
\text{and where }x\text{ or }y\text{ is an integer,}
\end{gathered}
\]
and put $\varepsilon=\pm1$ with the sign of the imaginary part of
\[
\frac{im}{p-p^0}
\]

Special rules: first class, $x$ an integer
\[
\begin{aligned}
2b\eta&=+a-2x+a\omega\\
2by&=+d-2ax+d\omega\\
2bX&=+\beta+2Bx+\beta\omega\\
2bY&=+\alpha-2Ax+\alpha\omega
\end{aligned}
\]
\noindent $\varepsilon=+1$ when $b$ is positive, $x,[y]$ are even, or when one condition holds;\par
\noindent $\varepsilon=-1$ when none or two conditions hold.

Second class, $y$ an integer
\[
\begin{aligned}
2a\eta&=-b+2y-b\omega\\
2ax&=+d-2by+d\omega\\
2aX&=+\alpha-2By+\alpha\omega\\
2aY&=-\beta+2Ay-\beta\omega
\end{aligned}
\]
\noindent $\varepsilon=+1$ when $a$ is positive, $[x],y$ are even, or when one condition holds;\par
\noindent $\varepsilon=-1$ when none or two conditions hold.
